\subsection{Đề 4}
\Opensolutionfile{ans}[ans/ans-2-OTHK2-De4]
\caulc
\begin{ex}%[2D4H2-2]
	Cho hàm số $f(x)$ liên tục trên $\mathbb{R}$ thỏa mãn $\displaystyle\int_{0}^{2}f(x)\mathrm{d}x=4$; $\displaystyle\int_{1}^{2}f(x)\mathrm{d}x=3$. Giá trị của biểu thức $\displaystyle\int_{0}^{1}f(x)\mathrm{d}x$ bằng
	\choice{$7$}{\True$1$}{$12$}{$0{,}75$}
	\loigiai{
		Ta có
		$$\displaystyle\int_{2}^{1}f(x)\mathrm{d}x=-\displaystyle\int_{1}^{2}f(x)\mathrm{d}x=-3.$$
		Suy ra 
		$$\displaystyle\int_{0}^{1}f(x)\mathrm{d}x=\displaystyle\int_{0}^{2}f(x)\mathrm{d}x+\displaystyle\int_{2}^{1}f(x)\mathrm{d}x=4-3=1.$$
	}
	
\end{ex}
\begin{ex} %[2H5N2-3]
	Đường thẳng $d$ đi qua $M\left( 2;0;-1 \right)$ và có véc tơ chỉ phương $\overrightarrow{u}=\left( 4;-6;2 \right)$ có phương trình
	\choice 
	{ \True $\heva{
			& x=2+2t \\ 
			& y=-3t \\ 
			& z=-1+t \\ 
		}$}
	{ $\heva{
			& x=-2+2t \\ 
			& y=-3t \\ 
			& z=1+t \\ 
		}$}
	{ $\heva{
			& x=-2+4t \\ 
			& y=-6t \\ 
			& z=1+2t \\ 
		}$}
	{ $\heva{
			& x=4+2t \\ 
			& y=-3t \\ 
			& z=2+t \\ 
		}$} 
	\loigiai{
		Véc-tơ $\overrightarrow{u}=\left( 4;-6;2 \right)$ là véc-tơ chỉ phương của $d$ nên véc-tơ $\overrightarrow{v}=(2;-3;1)$ cũng là là véc-tơ chỉ phương của $d$.\\
		Đường thẳng $d$ đi qua $M\left( 2;0;-1 \right)$ và có véc-tơ chỉ phương $\overrightarrow{v}=\left( 2;-3;1 \right)$ nên có phương trình là
		$$\heva{
			& x=2+2t \\ 
			& y=-3t \\ 
			& z=-1+t .\\ 
		}$$
	}
\end{ex} 
\begin{ex}%[2H5H1-6]
	Cho hai mặt phẳng $(P)\colon 2 x-y-z-3=0$ và $(Q)\colon x-z-2=0$. Góc giữa hai mặt phẳng $(P)$ và  $(Q)$ bằng
	\choice
	{$30^{\circ}$}
	{$45^{\circ}$}
	{\True $60^{\circ}$}
	{$90^{\circ}$}
	\loigiai{
		$(P)$ có vectơ pháp tuyến là $\overrightarrow{n}_1=(2 ; -1 ; -1)$.\\
		$(Q)$ có vectơ pháp tuyến là $\overrightarrow{n}_2=(1 ; 0 ; -1)$.\\
		Gọi $\varphi$ là góc giữa hai mặt phẳng $(P)$ và  $(Q)$ ta có $$\cos\varphi=\dfrac{\left| 2\cdot 1-1\cdot 0+ 1\cdot 1\right| }{\sqrt{2^2+(-1)^2+(-1)^2}\cdot\sqrt{1^2+0^2+(-1)^2}}=\dfrac{\sqrt{3}}{2}\Rightarrow \varphi=60^{\circ}.$$.
	}
\end{ex}
\begin{ex}%[2H5H2-6]
	Trong không gian $Oxyz$, tính góc giữa hai đường thẳng
	\begin{center}
		$d\colon \heva{&x=1+t\\&y=3-t\\&z=2t} \left(t \in \mathbb{R}\right)$ và $d'\colon \heva{&x=t'\\&y=1+2t'\\&z=3-t'} \left(t' \in \mathbb{R}\right)$
	\end{center}
	\choice	{$30^{\circ}$}
	{$45^{\circ}$}
	{\True $60^{\circ}$}
	{$90^{\circ}$}
	\loigiai
	{
		Đường thẳng $d$ và $d'$ lần lượt có các vectơ chỉ phương là $\overrightarrow{a}=\left(1;-1;2\right)$ và $\overrightarrow{a'}=\left(1;2;-1\right)$.\\
		Ta có
		$$\cos\left(d,d'\right)=\dfrac{|\overrightarrow{a}\cdot \overrightarrow{a'}|}{|\overrightarrow{a}|\cdot |\overrightarrow{a'}|}=\dfrac{|1\cdot 1-1\cdot 2+2\cdot (-1)|}{\sqrt{1^2+(-1)^2+2^2}\cdot \sqrt{1^2+2^2+(-1)^2}}=\dfrac{|-3|}{6}=\dfrac{1}{2}.$$
		Vậy $\left(d,d'\right)=60^{\circ}$.
	}
\end{ex}
\begin{ex}%[2H5H3-2]
	Cho mặt cầu $(S)\colon (x+1)^2+(y-2)^2+(z-1)^2=9$. 
	Toạ tâm $I$ và bán kính $R$ của $(S)$ là
	\choice
	{\True $I(-1 ; 2 ; 1)$ và $R=3$}
	{$I(1 ;-2 ;-1)$ và $R=3$}
	{$I(-1 ; 2 ; 1)$ và $R=9$}
	{$I(1 ;-2 ;-1)$ và $R=9$}
	\loigiai{
		Mặt cầu $(S)\colon (x+1)^2+(y-2)^2+(z-1)^2=9$
		có tâm $I(-1 ; 2 ; 1)$ và bán kính $R=3$.
	}
\end{ex}
\begin{ex}%[2H5H3-3]
	Mặt cầu tâm $I(-3 ; 0 ; 4)$ và đi qua điểm $A(-3 ; 0 ; 0)$ có phương trình là
	\choice
	{$(x-3)^2+y^2+(z+4)^2=4$}
	{$(x-3)^2-y^2+(z+4)^2=16$}
	{\True $(x+3)^2+y^2+(z-4)^2=16$}
	{$(x+3)^2+y^2+(z-4)^2=4$}
	\loigiai{
		Mặt cầu tâm $I(-3 ; 0 ; 4)$ và đi qua điểm $A(-3 ; 0 ; 0)$ nên bán kính $$R=IA=\sqrt{0+0+16}=4.$$
		Phương trình mặt cầu là $(x+3)^2+y^2+(z-4)^2=16$.
	}
\end{ex}
\begin{ex}%[2H5H3-3]
	Trong không gian với hệ tọa độ $Oxyz$, cho mặt cầu $(S)$ đi qua hai điểm $A(1;1;2)$, $B(3;0;1)$ và có tâm nằm trên trục $Ox$. Phương trình mặt cầu $(S)$ là
	\choice
	{${(x-1)}^2+y^2+z^2=\sqrt{5}$}
	{${(x+1)}^2+y^2+z^2=5$}
	{\True ${(x-1)}^2+y^2+z^2=5$}
	{${(x+1)}^2+y^2+z^2=\sqrt{5}$}
	\loigiai{
		Gọi $I$ là tâm mặt cầu $(S)$. Do $I\in Ox \Rightarrow I(x;0;0)$.\\
		Hai điểm $A(1;1;2)$, $B(3;0;1)$ thuộc mặt cầu $(S)$ nên
		\[IA=IB \Leftrightarrow IA^2=IB^2\Leftrightarrow (1-x)^2+1^2+2^2=(3-x)^2+0^2+1^2 \Leftrightarrow 4x=4 \Leftrightarrow x=1.\]
		Do đó $I(1;0;0)$ và $R^2=IA^2 ={(1-1)}^2+1^2+2^2=5$.\\
		Vậy phương trình mặt cầu $(S)$ là ${(x-1)}^2+y^2+z^2=5$.}
\end{ex}
\begin{ex}%[2H5V3-2]
	Trong không gian với hệ trục tọa độ $Oxyz$, cho ba điểm $A(1;0;0)$, $C(0;0;3)$, $B(0;2;0)$. Tập hợp các điểm $M$ thỏa mãn $MA^2=MB^2+MC^2$ là mặt cầu có bán kính là
	\choice
	{$R=2$}
	{\True $R=\sqrt{2}$}
	{$R=3$}
	{$R=\sqrt{3}$}
	\loigiai{Gọi $M(x;y;z)$ là điểm thỏa mãn bài toán. Ta có
		\begin{eqnarray*}
			MA^2=MB^2+MC^2&\Leftrightarrow&(x-1)^2+y^2+z^2=x^2+y^2+(z-3)^2+x^2+(y-2)^2+z^2\\
			&\Leftrightarrow&x^2+y^2+z^2+2x-4y-6z+12=0.
		\end{eqnarray*}
		Ta có $R=\sqrt{(-1)^2+2^2+3^2-12}=\sqrt{2}$.}
\end{ex}
\begin{ex}%[2D5N2-1]
	Cho $A$, $B$ là các biến cố của một phép thử $T$. Biết rằng $\mathrm{P}(B)>0$, xác suất của biến cố $A$ với điều kiện biến cố $B$ đã xảy ra được tính theo công thức nào sau đây?
	\choice
	{$\mathrm{P}(A|B)=\dfrac{\mathrm{P}(A)}{\mathrm{P}(B)}$}
	{$\mathrm{P}(A|B)=\dfrac{\mathrm{P}(A)}{\mathrm{P}(A B)}$}
	{\True $\mathrm{P}(A| B)=\dfrac{\mathrm{P}(A B)}{\mathrm{P}(B)}$}
	{$\mathrm{P}(A|B)=\dfrac{\mathrm{P}(A B)}{\mathrm{P}(A) \cdot \mathrm{P}(B)}$}
	\loigiai{ Dựa theo công thức tính xác suất biến cố $A$ với điều kiện $B$ thì $\mathrm{P}(A| B)=\dfrac{\mathrm{P}(A B)}{\mathrm{P}(B)}$ là đáp án đúng.
		
	}
\end{ex}

\begin{ex}%[2D5H2-1]
	Cho hai biến cố $A$ và $B$ có $\mathrm{P}(A)=0\text{,}8$; $\mathrm{P}(B)=0\text{,}5$ và $\mathrm{P}(A B)=0\text{,}2$. Xác suất của biến cố $A$ với điều kiện $B$ là
	\choice
	{\True$0\text{,}4$}
	{$0\text{,}5$}
	{$0\text{,}25$}
	{$0\text{,}625$}
	\loigiai{Ta có 
		$\mathrm{P}(A\mid B)=\dfrac{\mathrm{P}(A B)}{\mathrm{P}(B)}=\dfrac{0\text{,}2}{0\text{,}5}=0\text{,}4$.}
\end{ex}
\begin{ex}%[2D5H2-2]
	Cho hai biến cố $A$, $B$ với $\mathrm{P}(B)=0{,}6$; $\mathrm{P}(A|B) =0{,}7$ và $\mathrm{P}\left(A|\overline{B}\right)=0{,}4$. Khi đó, $\mathrm{P}(A)$ bằng
	\choice
	{$0{,}7$}
	{$0{,}4$}
	{\True $0{,}58$}
	{$0{,}52$}
	\loigiai{
		Ta có $\mathrm{P}\left(\overline{B}\right)= 1- \mathrm{P}(B) = 1-0{,}6 = 0{,}4$.\\
		Áp dụng công thức xác suất toàn phần, ta có
		\[\mathrm{P}(A) = \mathrm{P}(A|B)\cdot \mathrm{P}(B) + \mathrm{P}\left(A|\overline{B}\right)\cdot \mathrm{P}\left(\overline{B}\right)= 0{,}7\cdot 0{,}6 + 0{,}4\cdot 0{,}4=0{,}58.\]
	}
\end{ex}
\begin{ex}%[2D5H2-3] 
	\immini[thm]{Cho sơ đồ hình cây như hình bên. Biết $\mathrm{P}(A)=0{,}72$. Khi đó $\mathrm{P}(A \mid \overline{B})$ bằng
		\choice
		{\True $\dfrac{9}{16}$}
		{$\dfrac{16}{25}$}
		{$\dfrac{9}{25}$}
		{$\dfrac{3}{16}$}}
	{\begin{tikzpicture}[scale=.2,>=stealth]
			\tikzstyle{block} = [rectangle, draw, fill=blue!10\text{,} rounded corners, text centered, text width = 10em, minimum height = 2em]
			\node (c1) {};
			\node (c2)[above right = 1.5cm of c1] {$A$};
			%			\node at (0.5,5){\fbox{$0\text{,}7$}};
			%			\node at (0.5,-5){\fbox{$0\text{,}3$}};
			\node (c3) [below right= 1.5cm of c1]{$\overline{A}$};
			\node at (12,11.5){\fbox{$0\text{,}6$}};
			\node (c4) at (21.5, 12){$B$};
			\node (c5) at (21.5, 2){$\overline{B}$};
			\node at (12,3){\fbox{$0\text{,}4$}};
			\node (c6) at (21.5, -4){$B$};
			\node at (12,-4){\fbox{$0\text{,}2$}};
			\node (c7) at (21.5, -14){$\overline{B}$};
			\node at (12,-13){\fbox{$0\text{,}8$}};
			\draw[->] (c1.east) -- (c2.west);
			\draw[->] (c1.east) -- (c3.west);
			\draw[->] (c2.east) -- (c4.west);
			\draw[->] (c2.east) -- (c5.west);
			\draw[->] (c3.east) -- (c6.west);
			\draw[->] (c3.east) -- (c7.west);
	\end{tikzpicture}}
	\loigiai{
		Theo công thức Bayes ta có \begin{eqnarray*}
			\mathrm{P}\left( {A|\overline{B}} \right) &=& \dfrac{{\mathrm{P}\left( A \right) \mathrm{P}\left( {\overline{B}|A} \right)}}{{\mathrm{P}\left( \overline{B} \right)}} \\
			&=& \dfrac{{\mathrm{P}\left( A \right) \mathrm{P}\left( {\overline{B}|A} \right)}}{\mathrm{P}(A) \cdot \mathrm{P}(\overline{B}\mid A)+\mathrm{P}(\overline{A}) \cdot \mathrm{P}(\overline{B} \mid \overline{A})} \\
			&=& \dfrac{0{,}72 \cdot 0{,}4}{0{,}72 \cdot 0{,}4+0{,}28 \cdot 0{,}8} = \dfrac{9}{16}.
		\end{eqnarray*}
	}
\end{ex}
\Closesolutionfile{ans}
\indapan{6}{ans/ans-2-OTHK2-De4}

\cauds
\Opensolutionfile{ans}[ans/ans-2-OTHK2-De4-DS]
\begin{ex}%[2H5H2-4]
	Trong không gian $Oxyz$, cho hai đường thẳng $$\Delta_1\colon\dfrac{x-1}{-2}=\dfrac{y+2}{1}=\dfrac{z-4}{3},\,\Delta_2\colon \dfrac{x+1}{1}=\dfrac{y}{-1}=\dfrac{z+2}{3}.$$
	Xét các véc-tơ $\overrightarrow{u}_1=(-2;1;3)$ và $\overrightarrow{u}_2=(1;-1;3)$.
	\choiceTF
	{\True $\Delta_1$ đi qua điểm $M(1;-2;4)$ và có $\overrightarrow{u}_1$ là một véc-tơ chỉ phương}
	{\True $\Delta_2$ đi qua điểm $N(-1;0;-2)$ và có $\overrightarrow{u}_2$ là một véc-tơ chỉ phương}
	{$\left[ \overrightarrow{u}_1,\overrightarrow{u}_2\right]=(-6;-3;-1)$}
	{Hai đường thẳng $\Delta_1$ và $\Delta_2$ cắt nhau}
	\loigiai{
		\begin{itemchoice}
			\itemch  $\Delta_1$ đi qua điểm $M(1;-2;4)$ và có $\overrightarrow{u}_1$ là một véc-tơ chỉ phương.
			\itemch $\Delta_2$ đi qua điểm $N(-1;0;-2)$ và có $\overrightarrow{u}_2$ là một véc-tơ chỉ phương
			\itemch $\left[ \overrightarrow{u}_1,\overrightarrow{u}_2\right]=(6;3;1)$.
			\itemch Ta có $\overrightarrow{NM}=(2;-2;6)$.\\
			Vì $\left[ \overrightarrow{u}_1,\overrightarrow{u}_2\right]\cdot\overrightarrow{NM}=6\cdot 2+3\cdot (-2)+1\cdot 6=12\ne 0$ nên $\Delta_1$ và $\Delta_2$ chéo nhau.
		\end{itemchoice}
	}
\end{ex}


\begin{ex}%[2H5N2-7]
Trong không gian $Oxyz$, cho hai mặt phẳng $(P)\colon 2x+y+5z-3=0$ và \break $(Q)\colon 3x-y+z-2025=0$. Xét các véc-tơ $\overrightarrow{n}_1=(2;1;5)$ và $\overrightarrow{n}_2=(3;-1;1)$.
\choiceTF
{\True Mặt phẳng $(P)$ đi qua điểm $M(0;3;0)$ và có $\overrightarrow{n}_1$ là một véc-tơ pháp tuyến}
{Mặt phẳng $(Q)$ đi qua điểm $N(0;0;-2025)$ và có $\overrightarrow{n}_2$ là một véc-tơ pháp tuyến}
{$\overrightarrow{n}_1\cdot \overrightarrow{n}_2=1$}
{Góc giữa hai mặt phẳng $(P)$ và $(Q)$ bằng $30^\circ$}
\loigiai{
	Ta có $\overrightarrow{n}_1=(2;1;5)$ và $\overrightarrow{n}_2=(3;-1;1)$ lần lượt là véc-tơ pháp tuyến của  mặt phẳng $(P)$ và $(Q)$.
	\begin{itemchoice}
	\itemch Vì $2\cdot 0+3+5\cdot 0=0$ nên $M\in (P)$.
	\itemch Vì $3\cdot 0-0-2025-2025=-4050\ne 0$ nên $N\notin (Q)$.
	\itemch $\overrightarrow{n}_1\cdot \overrightarrow{n}_2=2\cdot 3+1\cdot (-1)+5\cdot 1=0\ne 1$.
	\itemch $\overrightarrow{n}_1\cdot \overrightarrow{n}_2=2\cdot 3+1\cdot (-1)+5\cdot 1=0$ suy ra $(P)$ và $(Q)$ vuông góc hay $((P),(Q))=90^\circ$.
\end{itemchoice}
}
\end{ex}

\begin{ex}%[2H5V3-2]
Trong không gian $Oxyz$, cho mặt cầu $(S)\colon (x-1)^2+(y-2)^2+(z-3)^2=1$ và điểm $A(2;3;4)$. Gọi $M$ là điểm thuộc $(S)$ sao cho đường thẳng $AM$ tiếp xúc với $(S)$.
\choiceTF
{\True Tâm mặt cầu $(S)$ là $I (1;2;3)$}
{ Điểm $A$ nằm trên mặt cầu $(S)$}
{Mặt phẳng $(P)\colon x+y+z-5=0$ cắt mặt cầu $(S)$ theo giao tuyến là một đường tròn có bán kính bằng $\dfrac{1}{\sqrt{3}}$}
{\True $M$ thuộc mặt phẳng có phương trình là $x+y+z-7=0$}
\loigiai{
Mặt cầu $(S)$ có tâm $I (1;2;3)$ và bán kính $R=1$.
\begin{itemchoice}
	\itemch Tâm mặt cầu $(S)$ là $I (1;2;3)$.
	\itemch Vì $AI=\sqrt{(2-1)^2+(3-2)^2+(4-3)^2}=\sqrt{3}>R$ nên $A$ nằm ngoài mặt cầu $(S)$.
	\itemch Ta có $\mathrm{d}(I,(P))=\dfrac{\left| 1+2+3-5\right| }{\sqrt{1^2+1^2+1^2}}=\dfrac{1}{\sqrt{3}}$.\\
	Mặt phẳng $(P)\colon x+y+z-5=0$ cắt mặt cầu $(S)$ theo giao tuyến là một đường tròn có bán kính $r=\sqrt{R^2-\mathrm{d}^2(I,(P))}=\dfrac{\sqrt{2}}{\sqrt{3}}$.
	\itemch Giả sử $M(x;y;z)$. Ta có $\overrightarrow{AM}=(x-2;y-3;z-4)$, $\overrightarrow{IM}=(x-1;y-2;z-3)$. Lại có $M\in (S)$ suy ra $(x-1)^2+(y-2)^2+(z-3)^2=1$ (*).\\
	Vì $AM$ tiếp xúc với $(S)$ nên $$\begin{aligned}
	AM\perp IM &\Leftrightarrow \overrightarrow{AM}\cdot \overrightarrow{IM}=0\\
	&\Leftrightarrow (x-2)(x-1)+(y-3)(y-2)+(z-4)(z-3)=0\\
	&\Leftrightarrow (x-1)^2+(y-2)^2+(z-3)^2 -(x-1)-(y-2)-(z-3)=0.\\
	&\Leftrightarrow 1-(x+y+z-6)=0~ (\text{Do} (*))\\
	&\Leftrightarrow x+y+z-7=0.
	\end{aligned}$$
\end{itemchoice}
}
\end{ex}


\begin{ex}%[2D5N2-4]
Cho hai biến cố $A$ và $B$ có $\mathrm{P}(B)=0{,}85$; $\mathrm{P}(A|B)=0{,}7$ và $\mathrm{P}(A|\overline{B})=0{,}5$.
\choiceTF
{$\mathrm{P}(\overline{B})=0{,}25$}
{\True$\mathrm{P}(A)=0{,}67$}
{\True $\mathrm{P}(B|A)=\dfrac{119}{134}$}
{ $\mathrm{P}(\overline{A})=0{,}3$}
\loigiai{
\begin{itemchoice}
	\itemch $\mathrm{P}(\overline{B})=1- 0{,}85=0{,}15$.
		\itemch $\mathrm{P}(A)=\mathrm{P}(B)\cdot \mathrm{P}(A|B)+\mathrm{P}(\overline{B})\cdot\mathrm{P}(A|\overline{B})=0{,}85\cdot 0{,}7+0{,}15\cdot 0{,}5=0{,}67$.
	\itemch $\mathrm{P}(B|A)=\dfrac{\mathrm{P}(B)\cdot\mathrm{P}(A|B) }{ \mathrm{P}(A)}=\dfrac{0{,}85\cdot 0{,}7}{0{,}67}=\dfrac{119}{134}$.
	\itemch $\mathrm{P}(\overline{A})=1-\mathrm{P}(A)=1-0{,}67=0{,}33$.
\end{itemchoice}
}
\end{ex}
\Closesolutionfile{ans}
\indapan{3}{ans/ans-2-OTHK2-De4-DS}

\Opensolutionfile{ans}[ans/ans-2-OTHK2-De4-KQ]
\caukq
\begin{ex}%[2H5H1-3]
	Cho đường thẳng $\Delta$ có phương trình $\Delta\colon \heva{& x=4+2t\\ & y=2-3t \\ & z=3+t}.$ Mặt phẳng $(P)$ đi qua điểm $M(2;-3;1)$ và chứa $\Delta$ có phương trình $-11z+by+cz+d=0$. Tính $b+c+d$.
	\shortans{$14$}
	\loigiai{
		Ta có $A(4;2;3)\in\Delta$ suy ra $\overrightarrow{AM}=(-2;-5;-2)$ và một véc-tơ chỉ phương của $\Delta$ là $\overrightarrow{u}_\Delta=(2;-3;1)$.\\
		Suy ra một véc-tơ pháp tuyến của $(P)$ là $\overrightarrow{n}_P=\left[\overrightarrow{AM},\overrightarrow{u}_\Delta\right]=(-11;-2;16)$.\\
		Và $(P)$ đi qua $M$ nên $(P)\colon-11(x-2)-2(y+3)+16(z-1)=0\Leftrightarrow -11x-2y+16z=0$.\\
		Từ đó suy ra $b=-2$, $c=16$, $d=0$.\\
		Vậy $b+c+d=-2+16+0=14$.}
\end{ex}
\begin{ex}%[2H5N2-7]
	Trong không gian với hệ tọa độ $Oxyz$, cho đường thẳng $d_1\colon\heva{&x=t\\&y=1-2t\\&z=-3t} (t\in\mathbb{R}) $ và đường thẳng $d_2\colon \dfrac{x}{-4}=\dfrac{y-1}{1}=\dfrac{z+1}{5}$. Tính góc giữa hai đường thẳng $d_1$, $d_2$.
	\shortans{$30^\circ$}
	\loigiai{
		$\overrightarrow{u}_1=(1;-2;-3)$ là một véc-tơ chỉ phương của $d_1$.\\
		$\overrightarrow{u}_2=(-4;1;5)$ là một véc-tơ chỉ phương của $d_2$.\\
		Ta có $\cos (d_1,d_2)=\dfrac{\vert 1\cdot (-4)+(-2)\cdot 1+(-3)\cdot 5\vert}{\sqrt{1^2+(-2)^2+(-3)^2}\cdot \sqrt{(-4)^2+1^2+5^2}}=\dfrac{\sqrt{3}}{2}$.\\
		Suy ra $(d_1,d_2)=30^\circ$.
	}
\end{ex}
\begin{ex}%[2H5N3-2]
	Trong không gian $Oxyz$, cho điểm $A(5;-1;3)$ và mặt cầu $(S)$ có phương trình $\left(x-1\right)^2+(y+1)^2+z^2=9$. Gọi $I$ là tâm mặt cầu $(S)$. Tính độ dài $IA$.
	\shortans{$5$}
	\loigiai{
		Phương trình mặt cầu $\left(x-1\right)^2+(y+1)^2+z^2=9$ có tâm $I \left(1; -1; 0\right)$.\\
		Ta có $\overrightarrow{IA}=(4;0;3)\Rightarrow IA=\sqrt{4^2+0^2+(-3)^2}=5$.
	}
\end{ex}
\begin{ex}%[2H5C3-2]
	Trong không gian với hệ trục tọa độ $Oxyz$, cho điểm $I(2;-1;1)$ và đường thẳng $d\colon \dfrac{x-2}{2}=\dfrac{y-1}{2}=\dfrac{z+1}{-1}$. Mặt cầu $(S)$ tâm $I$, bán kính $R$ cắt đường thẳng $d$  tại hai điểm $A$, $B$ sao cho $\triangle IAB$ vuông tại $I$. Tính $R^2$.
	\shortans{$8$}
	\loigiai{
		Mặt cầu tâm $I$ cắt đường thẳng $d$ tại $A,\, B$ nên  $A,\, B\in d$.\\ Khi đó $A(2+2t;1+2t;-1-t),\, B(2+2u;1+2u;-1-u)$ với $t\ne u$.\\
		Suy ra $\overrightarrow{IA}=(2t;2+2t;-2-t),\, \overrightarrow{IB}=(2u;2+2u;-2-u)$.\\
		Theo giả thiết ta có
		\begin{align*}
			& \heva{&\overrightarrow{IA}\cdot\overrightarrow{IB}=0\,(\text{do}\,\triangle IAB\,\text{vuông})\\&IA=IB\,(\text{do}\, I\, \text{là tâm hình cầu})} \Leftrightarrow \heva{&4tu+(2+2t)(2+2u)+(2+t)(2+u)=0\\&4t^2+(2+2t)^2+(2+t)^2=4u^2+(2+2u)^2+(2+u^2)}\\
			\Leftrightarrow\,& \heva{&9tu+6u+6t+8=0\\&9t^2+12t=9u^2+12u}
			\Leftrightarrow  \heva{&9tu+6u+6t+8=0\\&(t-u)(9t+9u+12)=0}
			\Leftrightarrow \,  \heva{&9tu+6u+6t+8=0\\& t=-u-\dfrac{4}{3}\,(\text{vì}\, t\ne u)}\\
			\Leftrightarrow\,& \heva{&9u\left(-u-\dfrac{4}{3}\right)+6\left(-u-\dfrac{4}{3}\right)+6u+8=0\\&t=-u-\dfrac{4}{3}}
			\Leftrightarrow \,  \heva{&-9u^2-12u=0\\&t=-u-\dfrac{4}{3}}
			\Leftrightarrow\, \hoac{&u=0;\, t=\dfrac{-4}{3}\\&u=-\dfrac{4}{3};\,t=0.}
		\end{align*}
		Nhận xét $A,\,B$ có vai trò như nhau nên ta chỉ cần xét trường hợp $u=0;\, t=-\dfrac{4}{3}$. Khi đó $B(2;1;-1),\, \overrightarrow{IB}=(0;2;-2)$.\\
		Suy ra bán kính hình cầu là $R=IB=2\sqrt{2}$.\\
		Suy ra $R^2=8$.
	}
\end{ex}
\begin{ex}%[2D5V2-2]
	Có hai chuồng thỏ. Chuồng I có 5 con thỏ đen và 10 con thỏ trắng. Chuồng II có 7 con thỏ đen và 3 con thỏ trắng. Trước tiên, từ chuồng II lấy ra ngẫu nhiên 1 con thỏ rồi cho vào chuồng I. Sau đó, từ chuồng I lấy ra ngẫu nhiên 1 con thỏ. Tính xác suất để con thỏ được lấy ra là con thỏ trắng. (Làm  tròn đến hàng phần trăm) 	\shortans{$0{,}64$}
	\loigiai{Xét biến cố $A$: \lq \lq Con thỏ được lấy ra từ chuồng II để cho vào chuồng I là con thỏ trắng\rq \rq~. 	Xét biến cố $B$: \lq \lq Con thỏ được lấy ra từ chuồng I là con thỏ trắng\rq \rq~. \\		
		Ta có $\mathrm{P}(B)=\mathrm{P}(A)\cdot \mathrm{P}(B \mid A)+\mathrm{P}(\overline{A}) \cdot \mathrm{P}(B \mid \overline{A})$.
		\begin{itemize}
			\item Tính $\mathrm{P}(A)$: Đây là xác suất để lấy ra ngẫu nhiên 1 con thỏ trắng từ chuồng II rồi cho vào chuồng I. Có $n\left(\Omega\right)=\mathrm{C}^1_{10}$, $n\left(A\right)=\mathrm{C}^1_3$. Vậy $\mathrm{P}(A)=\dfrac{3}{10}$.
			\item Tính $\mathrm{P}(\overline{A})$: $\mathrm{P}(\overline{A})=1-\mathrm{P}(A)=\dfrac{7}{10}$.
			\item Tính $\mathrm{P}(B\mid A)$: Đây là xác suất để lấy ra ngẫu nhiên 1 con thỏ trắng từ chuồng I với điều kiện đã chọn ra 1 con thỏ trắng từ chuồng II rồi cho vào chuồng I, tức là có 5 con thỏ đen và 11 con thỏ trắng ở trong chuồng I. Tương tự như trên ta có $\mathrm{P}(B\mid A)=\dfrac{11}{16}$.
			\item Tính $\mathrm{P}(B\mid \overline{A})$: Đây là để lấy ra ngẫu nhiên 1 con thỏ trắng từ chuồng I với điều kiện đã chọn ra 1 con thỏ đen từ chuồng II rồi cho vào chuồng I, tức là có 6 con thỏ đen và 10 con thỏ trắng ở trong chuồng I. Tương tự như trên  ta có $\mathrm{P}(B\mid \overline{A})=\dfrac{10}{16}$.
		\end{itemize}
		Vậy $\mathrm{P}(B)=\mathrm{P}(A)\cdot \mathrm{P}(B \mid A)+\mathrm{P}(\overline{A}) \cdot \mathrm{P}(B \mid \overline{A})=\dfrac{3}{10}\cdot \dfrac{11}{16}+\dfrac{7}{10}\cdot \dfrac{10}{16}=\dfrac{103}{160}\approx0{,}64$. \\
		Vậy xác suất để con thỏ được lấy ra là con thỏ trắng là $0{,}64$.
	}
\end{ex}
\begin{ex}%[2D5C2-3] 
	Ở một khu rừng nọ có $7$ chú lùn, trong đó có $4$ chú luôn nói thật, $3$ chú còn lại nói thật với xác suất $0{,}5$. Bạn Tuyết gặp ngẫu nhiên một chú lùn. Gọi $A$ là biến cố \lq\lq Chú lùn đó luôn nói thật\rq\rq\, và $B$ là biến cố \lq\lq Chú lùn đó tự nhận mình luôn nói thật\rq\rq. Biết rằng chú lùn mà bạn Tuyết gặp tự nhận mình là người luôn nói thật. Tính xác suất để chú lùn đó luôn nói thật. (Làm tròn hai chữ số thập phân)
	\shortans{$0{,}73$}
	\loigiai{
		Ta có		
		$\mathrm{P}(A)=\dfrac{4}{7}$; $\mathrm{P}(\overline{A})=\dfrac{3}{7}$; $\mathrm{P}(B\mid A)=1$; $\mathrm{P}(B\mid \overline{A})=0,5$.\\		
		Theo công thức Bayes, ta có
		\begin{eqnarray*}
			\mathrm{P}(B)&=&\mathrm{P}(A)\cdot \mathrm{P}(B|A)+\mathrm{P}(\overline{A})\cdot \mathrm{P}(B|\overline{A})\\		
			&=&\dfrac{4}{7}\cdot 1+\dfrac{3}{7}\cdot 0,5=\dfrac{11}{14}.
		\end{eqnarray*}
		Khi đó $$\mathrm{P}(A|B)=\dfrac{\mathrm{P}(AB)}{\mathrm{P}(B)}=\dfrac{\mathrm{P}(A)\cdot \mathrm{P}(B|A)}{\mathrm{P}(B)}=\dfrac{\dfrac{4}{7}\cdot 1}{\dfrac{11}{14}}=\dfrac{8}{11}\approx 0,73.$$
	}
\end{ex}

\Closesolutionfile{ans}
\indapan{6}{ans/ans-2-OTHK2-De4-KQ}
